次の3つの\bookterm{recurrence}{漸化式}を見てほしい.
\[
  a_{n+1}=a_n^2,
  \quad
  a_{n+1}=a_n^2-1,
  \quad
  a_{n+1}=a_n^2-2
\]
1つ目は,更新するたびに指数が2倍になる.
繰り返し代入すれば
\[
  a_n=a_1^{2^{n-1}}
\]
と分かる.$a_1>0$なら\bookterm{logarithm}{対数}で\bookterm{geometric}{等比数列}にもできる.
$a_1\le0$では実数の対数を最初から取れないが,元の更新は定義される.

2つ目は$-1$が対数の利用を妨げる.\bookterm{fixed}{不動点}は
\[
  \alpha=\alpha^2-1
  \quad\Longleftrightarrow\quad
  \alpha=\frac{1\pm\sqrt5}{2}
\]
であり,初項をこれに選べば定数解となる(\sectionref{節：初項依存型}).
ただし,$b_n=a_n-\alpha$と置いても,
\[
  b_{n+1}=b_n^2+2\alpha b_n
\]
と二次の項が残り,\bookterm{linear}{線形}にはならない.
ここまでの手法では閉じた\bookterm{general}{一般項}を得られないが,他の解法まで否定したわけではない.\footnote{ここでいう閉じた式の意味と,「解く」ことの範囲については\sectionref{節：漸化式の分類}で考える.}

3つ目の$-2$は,ある恒等式にちょうど合う.
基本型への変形が見つからなくても,\textbf{式に固有の構造}から解法が開けることがある.
この節では,その一致を見ていこう.
